\chapter[\emj{🌉}~MST: Kruskal y Prim (Árbol de Expansión Mínima)]{\emj{🌉}~MST: Kruskal y Prim (Árbol de Expansión Mínima)}

\begin{dsaquees}

Conectar todas las casas 🏠 de un pueblo con cable de internet gastando
lo MENOS posible. No importa la ruta exacta; solo que TODAS queden
conectadas y el cable total sea el más barato.

Dos estrategias para lograrlo:
\begin{itemize}
\item Kruskal: ordena todos los cables del más barato al más caro y ve
      agregándolos, saltándote cualquiera que forme un círculo inútil.
\item Prim: empieza en una casa y va "creciendo" la red, agregando
      siempre el cable más barato que toque una casa nueva.
\end{itemize}

Ambas llegan al mismo costo mínimo por caminos distintos.

\end{dsaquees}

\begin{dsaotraforma}

Un Árbol de Expansión Mínima (Minimum Spanning Tree) de un grafo no
dirigido y ponderado es un subconjunto de aristas que conecta TODOS los
vértices, sin ciclos, con el peso total mínimo. Tiene exactamente
$V-1$ aristas.

\textbf{Kruskal} es "goloso por aristas": ordena las aristas por peso y
las agrega si no crean ciclo. Para saber si crean ciclo usa Union-Find
(DSU). Ideal para grafos dispersos.

\textbf{Prim} es "goloso por nodos": crece un solo árbol desde un nodo,
eligiendo con una min-heap la arista más barata hacia un nodo fuera del
árbol. Ideal para grafos densos.

\end{dsaotraforma}

\begin{dsadiagrama}
\begin{verbatim}
Grafo no dirigido con pesos:

     1        3
  A ─── B ──────── C
  │  ╲2     ╱4   ╱
  │5  ╲   ╱   ╱6
  D ─── E ─
     7

Aristas ordenadas: A-B(1), A-E(2), B-C(3), B-E(4), A-D(5)...

Kruskal agrega (si no forma ciclo):
  A-B(1)  ✓   A-E(2)  ✓   B-C(3)  ✓
  B-E(4)  ✗ (A,B,E ya conectados -> ciclo)
  A-D(5)  ✓
MST = {A-B, A-E, B-C, A-D}, peso total = 1+2+3+5 = 11
(V=5 nodos -> V-1 = 4 aristas)
\end{verbatim}
\end{dsadiagrama}

\begin{lstlisting}[language=C++]
KRUSKAL (goloso por aristas + Union-Find)
1. Ordena todas las aristas por peso ascendente.
2. Recorre las aristas: si u y v están en componentes distintas
   (find(u) != find(v)), agrégala al MST y hazles union().
3. Para cuando el MST tenga V-1 aristas.

📝 PSEUDOCÓDIGO (Kruskal)
──────────────────────────────────────────────────────────────

función kruskal(V, aristas):
    ordenar aristas por peso
    dsu = UnionFind(V)
    mst = 0; usadas = 0
    PARA cada (w, u, v) en aristas:
        SI dsu.find(u) != dsu.find(v):
            dsu.union(u, v)
            mst += w; usadas++
            SI usadas == V-1: romper
    return mst

💻 CÓDIGO C++ (Kruskal)
──────────────────────────────────────────────────────────────

#include <iostream>
#include <vector>
#include <algorithm>
using namespace std;

struct DSU {
    vector<int> p, r;
    DSU(int n): p(n), r(n, 0) { for (int i=0;i<n;i++) p[i]=i; }
    int find(int x){ return p[x]==x ? x : p[x]=find(p[x]); }
    bool unite(int a,int b){
        a=find(a); b=find(b);
        if(a==b) return false;
        if(r[a]<r[b]) swap(a,b);
        p[b]=a; if(r[a]==r[b]) r[a]++;
        return true;
    }
};

// aristas: {peso, u, v}
long long kruskal(int V, vector<array<int,3>>& edges) {
    sort(edges.begin(), edges.end()); // ordena por peso (primer campo)
    DSU dsu(V);
    long long total = 0; int usadas = 0;
    for (auto& e : edges) {
        if (dsu.unite(e[1], e[2])) {   // no forma ciclo
            total += e[0];
            if (++usadas == V - 1) break;
        }
    }
    return total;
}

💻 CÓDIGO C++ (Prim con min-heap)
──────────────────────────────────────────────────────────────

#include <queue>
long long prim(int V, vector<vector<pair<int,int>>>& adj) {
    // adj[u] = lista de {vecino, peso}
    vector<bool> inMST(V, false);
    priority_queue<pair<int,int>, vector<pair<int,int>>,
                   greater<>> pq;    // min-heap {peso, nodo}
    pq.push({0, 0});
    long long total = 0;
    while (!pq.empty()) {
        auto [w, u] = pq.top(); pq.pop();
        if (inMST[u]) continue;
        inMST[u] = true; total += w;
        for (auto [v, pw] : adj[u])
            if (!inMST[v]) pq.push({pw, v});
    }
    return total;
}

⏱️ COMPLEJIDAD (Big-O)
──────────────────────────────────────────────────────────────

  Algoritmo   │ Tiempo          │ Mejor para
  ────────────┼─────────────────┼────────────────
  Kruskal     │ O(E log E)      │ Grafos dispersos
  Prim (heap) │ O(E log V)      │ Grafos densos
\end{lstlisting}

\begin{dsaentrevistas}

✅ Kruskal reusa Union-Find: si ya dominas DSU, Kruskal es casi gratis.
   Es la razón por la que se enseñan juntos.

💡 "Min Cost to Connect All Points" (LeetCode 1584): MST clásico sobre
   distancias de Manhattan. Prim suele ganar por ser el grafo denso.

⚠️ MST solo en grafos NO dirigidos y conexos. Si el grafo está
   desconectado no existe MST (existe un "bosque" de expansión mínima).

🔑 Corte (cut property): la arista más barata que cruza cualquier
   partición del grafo SIEMPRE está en algún MST. Es la justificación
   teórica de por qué el enfoque goloso funciona aquí.

\end{dsaentrevistas}

\begin{dsaresumen}

• MST: conecta todos los nodos, sin ciclos, con peso mínimo (V-1 aristas).
• Kruskal: ordena aristas + Union-Find; O(E log E); grafos dispersos.
• Prim: crece un árbol con min-heap; O(E log V); grafos densos.
• Ambos son golosos y dan el mismo costo mínimo.
• Solo en grafos no dirigidos y conexos.
• Kruskal es la mejor excusa para practicar Union-Find.
════════════════════════════════════════════════════════════════

\end{dsaresumen}
